\chapter{Contraceptive Implant Insertion}

\section*{Overview}
Contraceptive implant insertion places a small, flexible rod (typically etonogestrel) under the skin of the inner upper arm. It provides long-acting, reversible contraception lasting up to 5 years.

\section*{Alternative Names}
Etonogestrel implant placement, Implanon/Nexplanon insertion, subcutaneous contraceptive implant placement

\section*{Principle}
The implant releases a steady dose of the progestin etonogestrel into the circulation. Progestin suppresses ovulation, thickens cervical mucus to impede sperm transport, and thins the endometrium. Contraceptive effect is established within 24 hours if inserted during the first 5 days of menses.

\section*{History}
The first contraceptive implant, a levonorgestrel system, was approved in 1991. The single-rod etonogestrel implant replaced earlier multi-rod devices and was approved in the United States in 2006.

\section*{Why It Is Done}
Patients request contraceptive implant insertion for highly effective, long-acting reversible contraception without daily adherence. It is also used in appropriate candidates as an alternative to other hormonal methods and may help manage heavy or painful menstruation.

\section*{Is This a Primary or Secondary Test or Procedure}
It is a primary contraceptive procedure providing long-acting reversible contraception.

\section*{How It Is Performed}
The clinician numbs the inner upper arm with local anesthetic, makes a small incision (2 mm) with the proprietary applicator, and inserts the rod just under the skin. Placement is confirmed by palpating the rod and, when indicated, by ultrasound or radiography. No sutures are typically required; a pressure dressing and bandage are applied.

\section*{Preparation}
A pregnancy test is performed to exclude pregnancy before insertion. Timing should be within the first 5 days of the menstrual cycle or per product labeling. No fasting or other specific preparation is required.

\section*{Contraindications and Precautions}
Contraindications include known or suspected pregnancy, active thromboembolic disease, and hypersensitivity to the implant components. Caution is warranted in women with unexplained vaginal bleeding, liver disease, or a history of venous thromboembolism.

\section*{Risks}
Risks include insertion-site pain, bruising, infection, and scarring. Implant migration, breakage, and difficult removal are rare.

\section*{Aftercare and Recovery}
The bandage is removed after 24 hours and the insertion site kept clean and dry. Irregular bleeding is common in the first months; fertility returns promptly after removal.

\section*{Outcome and Follow-up}
The implant is considered effective if it remains palpable subcutaneously and the patient has no new bleeding pattern concerns. Removal is required if the rod is no longer palpable or if contraception is no longer desired.

\section*{Expected Results}
Successful placement of a single, palpable implant in the subcutaneous tissue of the upper arm with no signs of infection or migration.

\section*{Follow-up}
Follow-up at 4-6 weeks is offered to confirm the site is healing and to address bleeding concerns. The implant should be removed and replaced by 5 years.

\section*{When to Seek Medical Care}
Follow the procedure team's specific instructions. Seek urgent medical assessment for severe or worsening pain, uncontrolled bleeding, fever or signs of infection, trouble breathing, chest pain, fainting, new weakness or confusion, inability to urinate, or any rapidly worsening symptom. The appropriate warning signs depend on the procedure and the patient's medical condition.

\section*{References}
\begin{enumerate}
  \item MedlinePlus. Contraceptive Implant Insertion. U.S. National Library of Medicine; 2025.
  \item Current Medical Diagnosis and Treatment. McGraw-Hill; 2025.
\end{enumerate}

\clearpage
